% Auxiliary lemma for the direct integral cycle detector; unregistered.
\paragraph{Auxiliary lemma (proved and independently audited):
divisibility of the relevant $K_2$-cohomology.}
Let $k$ be an algebraically closed field of arbitrary characteristic,
let $C/k$ be a smooth connected projective curve, and let $W/k$ be a
smooth cubic surface. For every $m\geq0$, the groups
\[
 H^1(\mathbb P^m\times W\times C,\mathcal K^M_2),
 \qquad
 H^1(\mathbb P^m\times W\times C,\mathcal K^Q_2)
\]
are divisible abelian groups. Here $\mathcal K^M$ and $\mathcal K^Q$
denote the Zariski sheaves of Milnor and Quillen $K$-groups,
respectively. This concerns sheaf cohomology, and is not inferred from
a decomposition of the group $K_1(\mathbb P^m\times W\times C)$.

\begin{proof}
We write
\[
 G^{a,b}(T)=H^a(T,\mathcal K^M_b)
\]
for a smooth $k$-variety $T$. Its integral Gersten resolution has
degree-$j$ term
\[
 \bigoplus_{x\in T^{(j)}}K^M_{b-j}(k(x)).
\]
The exactness is \citet[Theorem 7.1, author pp.~21--22]{cyKerz2009Milnor}.
Consequently these are precisely the corresponding graded Chow groups
with coefficients in Milnor $K$-theory. We use the localization
sequence of \citet[Section 5, p.~356]{cyRost1996}, the identification
with sheaf cohomology in his Corollary 6.5, p.~360, affine homotopy
invariance in Proposition 8.6, p.~370, and the pullback construction
and open-restriction compatibility in Section 12, pp.~382--383.
All these statements are integral and valid in every characteristic.

First, if $S$ is smooth connected and projective and $d\geq1$, put
$T=\mathbb P^d\times S$, $H=\mathbb P^{d-1}\times S$, and
$U=\mathbb A^d\times S$. Localization gives the exact portion
\[
 G^{0,1}(H)\longrightarrow G^{1,2}(T)
 \longrightarrow G^{1,2}(U).
\]
The last group is canonically $G^{1,2}(S)$ by affine homotopy
invariance. Its elements lift to $T$ by the projection pullback.
Moreover
\[
 G^{0,1}(H)=\Gamma(H,\mathcal O_H^*)=k^*.
\]
It follows that the sum of projection pullback and hyperplane
pushforward is a surjection
\[
 G^{1,2}(S)\oplus k^*\twoheadrightarrow G^{1,2}(T).
 \tag{K2.1}
\]
We need this generation assertion, rather than a full projective-bundle
decomposition.

Second, let $\pi:\widetilde T\to T$ be the blow-up of a smooth
connected projective variety along a smooth connected centre of
codimension at least two, and let $j:E\hookrightarrow\widetilde T$
be the exceptional divisor. Suppose $E$ is connected, as it is for
the centres used below. The proper pushforward and the pullback to a
smooth target give, for $b\in G^{1,2}(\widetilde T)$, the class
\[
 b-\pi^*\pi_*b.
\]
Its restriction to $\widetilde T\setminus E$ is zero: on this open
set $\pi$ is an isomorphism, proper pushforward commutes with open
restriction, and the pullback construction is local on the target.
The former compatibility is \citet[Proposition 4.1(3), p.~351]{cyRost1996};
the latter is stated immediately before Theorem 12.1, p.~382.
Localization therefore expresses this difference as $j_*u$ with
\[
 u\in G^{0,1}(E)=k^*.
\]
Thus
\[
 G^{1,2}(T)\oplus k^*\twoheadrightarrow
 G^{1,2}(\widetilde T),\qquad (c,u)\longmapsto\pi^*c+j_*u.
 \tag{K2.2}
\]
In particular no assertion about an unverified blow-up formula for
sheaf cohomology is necessary. The proper pushforward here preserves
both indicated indices because source and target have the same
dimension; the exceptional pushforward raises both indices by one.

Over an algebraically closed field every smooth cubic surface is the
blow-up of $\mathbb P^2$ at six points. Apply (K2.1) to
$\mathbb P^2\times C$, apply (K2.2) to the six smooth centres
$\{q_i\}\times C$, and apply (K2.1) once more for $\mathbb P^m$
when $m\geq1$. We obtain a surjection
\[
 G^{1,2}(C)\oplus(k^*)^{7+\epsilon_m}
 \twoheadrightarrow G^{1,2}(\mathbb P^m\times W\times C),
 \qquad
 \epsilon_m=\begin{cases}0&m=0,\\1&m\geq1.\end{cases}
 \tag{K2.3}
\]
For the curve the same Gersten resolution gives
\[
 G^{1,2}(C)=\operatorname{coker}\left(
 K^M_2(k(C))\longrightarrow\bigoplus_{c\in C(k)}k^*\right).
\]
Since $k$ is algebraically closed, $k^*$ is divisible. Direct sums and
quotients of divisible groups are divisible, so the source and target
of (K2.3) are divisible.

Finally, the degree-two Gersten complexes for Milnor and Quillen
$K$-theory agree: their only nonzero terms involve the field groups
$K_2$, $K_1$, and $K_0$, and their residue maps agree. This is also
the degree-two case of \citet[Remark 5.4, p.~357]{cyRost1996}, together
with the smooth Gersten resolution recorded in his Theorem 6.1 and
Corollary 6.5. The canonical map therefore identifies the two first
cohomology groups in the assertion. This comparison is made on the
smooth variety; it asserts no Milnor--Quillen comparison theorem for
an arbitrary nonreduced scheme.
\end{proof}

\noindent\emph{Source scope.} The complete Rost primary PDF and his
author-page corrections were opened. The false, unused flatness remark
on p.~377 and the corrected signs in (14.3)--(14.4) play no role in
this proof. The cubic-surface blow-up fact is the same all-characteristic
input independently checked in C.3bi. No direct-sum decomposition is
asserted by this auxiliary lemma.

\paragraph{Higher Milnor-degree extension (proved and independently
audited).}
Let $W_1,\ldots,W_s$ be smooth cubic surfaces over the algebraically
closed field $k$, and let $C$ be as above. For every $m,s,a\geq0$,
\[
 H^a\bigl(\mathbb P^m\times W_1\times\cdots\times W_s\times C,
                 \mathcal K^M_{a+1}\bigr)
 \quad\hbox{is divisible}.
 \tag{K2.4}
\]
In particular this includes the group with $a=r-1$ and Milnor degree
$r$ that occurs before a codimension-$r$ relative cycle detector.
No analogous higher-degree Milnor--Quillen identification is asserted.

\begin{proof}
We say that a smooth connected projective variety $S$ has property
$\mathcal D$ if $G^{a,a+1}(S)$ is divisible for every $a\geq0$;
groups with negative first index are interpreted as zero. The curve
$C$ has this property: degree zero is $k^*$, degree one is the
cokernel displayed above, and all higher groups vanish because the
Gersten resolution has length one.

For $T=\mathbb P^d\times S$, the same localization and projection
argument gives, in every bidegree, a surjection
\[
 G^{a,b}(S)\oplus
 G^{a-1,b-1}(\mathbb P^{d-1}\times S)
 \twoheadrightarrow G^{a,b}(\mathbb P^d\times S).
 \tag{K2.5}
\]
For $a=0$ the second term is zero; the assertion follows directly from
the exact localization sequence and the lifting of every open class
by projection pullback. Induction on $d$ proves that $\mathcal D$
is preserved by product with any projective space.

Similarly, for a blow-up $\pi:\widetilde T\to T$ of smooth
varieties with exceptional divisor $E$, the same expression
$b-\pi^*\pi_*b$ and localization give
\[
 G^{a,b}(T)\oplus G^{a-1,b-1}(E)
 \twoheadrightarrow G^{a,b}(\widetilde T).
 \tag{K2.6}
\]
In the product blow-up of $\mathbb P^2\times S$ along
$\{q\}\times S$, the exceptional divisor is $\mathbb P^1\times S$.
If $S$ has $\mathcal D$, then both summands of (K2.6) with
$b=a+1$ are divisible, by the preceding projective-space argument.
After six such blow-ups we conclude that $W\times S$ has
$\mathcal D$ for every smooth cubic surface $W$. Repeating this
step for $W_1,\ldots,W_s$, and then adjoining $\mathbb P^m$, proves
(K2.4). Every operation preserves the total Milnor grading with the
displayed index shifts, so the argument never substitutes ordinary
$K$-groups for the required Milnor $K$-cohomology.
\end{proof}

\paragraph{Consequence for a proposed direct detector
(conditional only on the stated relative-sheaf data).}
Let $E$ be a square-zero thickening of the smooth variety $V$ above,
and suppose there is an exact sequence of Zariski sheaves
\[
 0\longrightarrow\mathcal J\longrightarrow\mathcal K_2(E)
 \longrightarrow\mathcal K_2(V)\longrightarrow0,
 \qquad p\mathcal J=0
\]
in characteristic $p$. Then
\[
 H^2(V,\mathcal J)\longrightarrow H^2(E,\mathcal K_2(E))
\]
is injective. Indeed, its kernel is the image of the connecting map
from the divisible group $H^1(V,\mathcal K_2(V))$ into a group killed
by $p$, and every such homomorphism is zero. If a product of line
classes has a relative lift whose image under a sheaf map
$\mathcal J\to L\otimes\Omega^1_{V/k}$ is nonzero in $H^2$, the
absolute product class is nonzero. The construction and compatibility
of that relative lift and sheaf map must be proved separately; no
such compatibility is inferred merely from divisibility.
